\begin{textsample}{POS Dim 3 – to – Score 18.00 – to/to\_inf\_28.txt}  \label{ex:f3_pos_093}
O campo de experiência “ Corpo, \textbf{Gestos} e Movimentos ” O campo de experiência “ Corpo, \textbf{Gestos} e Movimentos ” aborda as experiências corporais que são prioritárias e fundamentais para as \textbf{crianças}, pois, o corpo e o pensamento estão estreitamente relacionados na infância.
O corpo carrega consigo, não somente características físicas e biológicas, mas também marcas do pertencimento social que ecoam em quem se é e nas experiências pessoais, em relação ao gênero, à etnia, à classe social, à religião e à sexualidade.
Campo de experiência : corpo, \textbf{gestos} e movimentos As \textbf{crianças} \textbf{brincam} com seu corpo, se comunicam e se expressam, por meio das suas múltiplas linguagens.
Desde \textbf{bebês}, o corpo é instrumento essencial para aprendizagem, se relaciona constantemente com o estético e o \textbf{sensorial}, ele expressa, sente, e proporciona que o outro sinta também.
Esse campo de experiência permite que a \textbf{criança} conheça e reconheça suas sensações, funções corporais e, nos seus \textbf{gestos} e movimentos, identifique suas potencialidades e limites, desenvolvendo, ao mesmo tempo, a consciência corporal, reconhecendo o processo de diferenciação do eu, do outro e da construção de sua identidade.
As experiências motoras permitem integrar as diferentes linguagens e, por meio dos movimentos e dos \textbf{gestos}, a \textbf{criança} explora o meio em que está inserida, desenvolve significados sobre os objetos, as pessoas e o mundo ; descobre seus próprios limites, enfrenta novos desafios, conhece a si mesma, expressa seus sentimentos, organiza e se localiza espacialmente.
Também desenvolve outras habilidades que contribuem para o desenvolvimento emocional, intelectual, afetivo, social e cognitivo.
Nesse sentido, o movimento e a linguagem corporal estão relacionados ao desenvolvimento e à aprendizagem.
Então, percebe -se a importância de se compreender e perceber que o corpo e o movimento constituem uma linguagem, uma forma comunicativa do ser humano se relacionar com o mundo.
Por meio da linguagem corporal, é dada a oportunidade à \textbf{criança} de experimentar as manifestações culturais, se comunicar e se expressar com outras \textbf{crianças}, com o \textbf{adulto} e com o mundo.
O contato com diferentes parceiros, materiais e espaços possibilita às \textbf{crianças} investigar as possibilidades de movimento que eles oferecem.
Em função disto, os espaços e as atividades cotidianas na Educação \textbf{Infantil} devem ser estruturados para possibilitar que as \textbf{crianças} indígenas, as ribeirinhas, as do campo e as \textbf{crianças} dos centros urbanos, estas cada vez mais limitadas no ambiente \textbf{doméstico}, cotidianamente, explorem seus \textbf{gestos} e movimentos de forma \textbf{lúdica}, considerando ainda as necessidades específicas de \textbf{movimentação} dos \textbf{bebês} e também das \textbf{crianças} com necessidades educacionais especiais ( BRASIL/MEC, 2018, p. 38 ).
A interação e o \textbf{brincar}, nesse contexto, possibilitam que a \textbf{criança} se expresse de forma \textbf{lúdica}, interaja com os objetos, com os outros e com o mundo, construindo significados e consciência corporal que a possibilite explorar movimentos, \textbf{gestos}, \textbf{sons}, formas, \textbf{texturas}, emoções, para que possa se expressar de forma criativa, fazer descobertas, hipóteses e utilizar diversificadas formas de linguagem corporal, para conhecer -se e construir sua própria identidade.
Corpo, \textbf{gestos} e movimentos - CONVIVER com \textbf{crianças} e \textbf{adultos} e experimentar, de múltiplas formas, a gestualidade que marca sua cultura e está presente nos \textbf{cuidados} pessoais, dança, música, teatro, artes circenses, jogos, \textbf{escuta} de histórias e \textbf{brincadeiras}. - \textbf{BRINCAR}, utilizando movimentos para se expressar, explorar espaços, objetos e situações, \textbf{imitar}, jogar, imaginar, interagir e utilizar \textbf{criativamente} o repertório da cultura corporal e do movimento. - PARTICIPAR de diversas atividades de \textbf{cuidados} pessoais e do contexto social, de \textbf{brincadeiras}, encenações teatrais ou circenses, danças e músicas ; desenvolver práticas corporais e autonomia para cuidar de si, do outro e do ambiente. - EXPLORAR amplo repertório de movimentos, \textbf{gestos}, olhares, \textbf{sons} e \textbf{mímicas} ; descobrir modos de ocupação e de uso do espaço com o corpo e adquirir a compreensão do seu corpo no espaço, no tempo e no grupo. - EXPRESSAR corporalmente emoções, ideias e opiniões, tanto nas relações cotidianas como nas \textbf{brincadeiras}, dramatizações, danças, músicas, contação de histórias, dentre outras manifestações, empenhando -se em compreender o que outros também expressam. - CONHECER -SE nas diversas oportunidades de interações e explorações com seu corpo ; reconhecer e valorizar o seu pertencimento de gênero, étnico-racial e religioso.
Brasil, 2016 ( 2a Versão da BNCC )

% matched lemmas: adulto, bebê, brincadeira, brincar, criança, criativamente, cuidado, doméstico, escuta, gesto, imitar, infantil, lúdico, movimentação, mímico, sensorial, som, textura
\end{textsample}
